\centerline{\bf \Huge  ТЧ All Stars}
\title{ТЧ All Stars}

\begin{flushright}
    \scriptsize Вот и все!\\ Неплохая получилась история: \\интересная, весёлая, порой немного грустная, а главное поучительная. \\Она научила быть нас смелыми и не бояться вызовов, которые готовит нам жизнь. \\Помогала нам добиваться поставленных целей несмотря ни на что... \\Но, самое главное, что у этой истории счастливый конец.
\end{flushright}

\problem{1}Докажите, что среди любых n чисел можно выбрать несколько таким образом, что их сумма будет делиться на n.

\problem{2}\textbf{Теорема Вильсона(Алине сдавать нельзя)}. Докажите, что $(n-1)! \equiv -1 \ (mod \ n) \longleftrightarrow n - \text{просто число}$

\problem{3} Докажите, что $C_p^k$ всегда делится на $p$ при любом $0<k<p$

\problem{4}\textbf{Формула двоечника.} Докажите, что $(a+b)^p \equiv a^p + b^p \ (mod \ 2)$, где $p$ простое число.

\problem{5}\textbf{Различные вариации теоремы Евклида.}

а)Докажите, что существует бесконечно много простых чисел вида  $p = 4k + 1.$

б)Докажите, что существует бесконечно много простых чисел вида  $p = 4k + 3.$

в)Докажите, что существует бесконечно много простых чисел вида  $p = 6k + 5.$

\problem{6}\textbf{Первая классика суммы дробей.} Докажите, что число $1 + \dfrac{1}{2} + \dfrac{1}{3} + \dotsc+ \dfrac{1}{n}$  при  $n > 1$  не будет целым.

\problem{7}\textbf{Вторая классика суммы дробей.} Пусть $p>5$ простое число и $1+\dfrac{1}{2}+\dfrac{1}{3}+\dotsc+\dfrac{1}{p-1}=\dfrac{m}{n}.$

а) Докажите, что $m \ \vdots \ p$

б) Докажите, что $m \ \vdots \ p^2$

\problem{8}\textbf{Классика уравнений в целых числах.}

a) Докажите, что уравнение  $3x^2 + 2 = y^2$  нельзя решить в целых числах.

б)Решить в целых числах уравнение  (2x + y)(5x + 3y) = 7.

в)Найдите все натуральные m и n, для которых  $m! + 12 = n^2$.

г)Решите в натуральных числах уравнение  $x^y = y^x$  при  $x \neq y$.

\problem{9}\large{\textbf{ЗА ИСПОЛЬЗОВАНИЕ ЭТОГО ГРОБА ФИЗФАК МГУ ЗАБАНИЛ МЕНЯ НА ОЛИМПИАДЕ...}}

Докажите, что между $n$ и $2n$ всегда есть простое число, для любого n>1.







